% !TEX root = ../DoAn.tex
\documentclass[../DoAn.tex]{subfiles}
\begin{document}

Chương này trình bày toàn bộ kết quả phân tích và thiết kế hệ thống, từ việc xác định yêu
cầu, mô hình hóa kiến trúc tổng thể, thiết kế cơ sở dữ liệu, thiết kế API và giao tiếp,
đến thiết kế luồng sự kiện, bảo mật, máy trạng thái và các luồng nghiệp vụ chính. Đây là
chương có mật độ sơ đồ và bảng biểu cao nhất, phản ánh toàn bộ quá trình chuyển đổi từ yêu
cầu nghiệp vụ sang thiết kế kỹ thuật chi tiết. Các quyết định thiết kế trong chương này
dựa trên nền tảng lý thuyết đã trình bày ở Chương~\ref{chapter:Related_works} và được hiện
thực hóa bằng các công nghệ đã phân tích ở Chương~\ref{chapter:Methodology}.

% =======================================================================
\section{Phân tích yêu cầu hệ thống}
\label{section:4.1}

\textbf{Yêu cầu chức năng.} Hệ thống được xây dựng để đáp ứng tập hợp đầy đủ các chức năng
của một nền tảng thương mại điện tử, phục vụ hai nhóm người dùng chính là khách hàng và quản
trị viên. Bảng~\ref{table:functional-requirements} tổng hợp các yêu cầu chức năng được nhóm
theo miền nghiệp vụ.

{%
\small\setlength{\tabcolsep}{3pt}\renewcommand{\arraystretch}{1.15}%
\begin{longtable}{|p{3.1cm}|p{2.4cm}|p{8.8cm}|}
\caption{Yêu cầu chức năng theo miền nghiệp vụ}
\label{table:functional-requirements} \\
\hline
\textbf{Miền nghiệp vụ} & \textbf{Đối tượng} & \textbf{Yêu cầu chức năng} \\ \hline
\endfirsthead
\hline
\textbf{Miền nghiệp vụ} & \textbf{Đối tượng} & \textbf{Yêu cầu chức năng} \\ \hline
\endhead
Định danh \& Người dùng & Khách hàng & Đăng ký tài khoản, đăng nhập, làm mới token \\ \cline{3-3}
                        & Khách hàng & Quản lý thông tin cá nhân, quản lý địa chỉ giao hàng \\ \hline
Sản phẩm \& Danh mục    & Khách hàng & Xem danh sách sản phẩm, lọc theo danh mục/thương hiệu, xem chi tiết \\ \cline{3-3}
                        & Quản trị viên & Tạo, sửa, xóa sản phẩm, danh mục, thương hiệu \\ \hline
Tìm kiếm                & Khách hàng & Tìm kiếm toàn văn theo từ khóa, gợi ý sản phẩm liên quan \\ \hline
Giỏ hàng                & Khách hàng & Thêm, sửa số lượng, xóa sản phẩm; hỗ trợ giỏ hàng khách vãng lai \\ \hline
Đơn hàng                & Khách hàng & Tạo đơn hàng (COD/VNPAY), xem danh sách đơn hàng, xem chi tiết trạng thái \\ \cline{3-3}
                        & Hệ thống   & Tự động xử lý luồng Saga: trừ tồn kho, xác nhận thanh toán, hủy đơn quá hạn \\ \hline
Thanh toán              & Khách hàng & Thanh toán online qua VNPAY, xem kết quả thanh toán \\ \cline{3-3}
                        & Hệ thống   & Xử lý IPN callback từ VNPAY, tự động timeout thanh toán quá hạn \\ \hline
Khuyến mãi              & Khách hàng & Xem danh sách voucher khả dụng, áp dụng voucher khi đặt hàng \\ \cline{3-3}
                        & Quản trị viên & Tạo, sửa, xóa voucher; cấu hình điều kiện sử dụng \\ \hline
Flash Sale              & Khách hàng & Xem danh sách chiến dịch, mua hàng flash sale với số lượng giới hạn \\ \cline{3-3}
                        & Quản trị viên & Tạo, kích hoạt, kết thúc chiến dịch; cấu hình sản phẩm và số lượng \\ \hline
Đánh giá                & Khách hàng & Tạo, sửa, xóa đánh giá sản phẩm (chỉ sau khi mua hàng thành công) \\ \hline
Nội dung                & Quản trị viên & Quản lý banner, bài viết blog \\ \hline
Thông báo               & Hệ thống   & Gửi email xác nhận đơn hàng và hủy đơn \\ \hline
\end{longtable}
}

\textbf{Yêu cầu phi chức năng.} Bên cạnh các yêu cầu chức năng, hệ thống phải đáp ứng các yêu
cầu phi chức năng then chốt của một hệ thống phân tán, được tổng hợp trong Bảng~\ref{table:nfr}.

{%
\small\setlength{\tabcolsep}{3pt}\renewcommand{\arraystretch}{1.15}%
\begin{longtable}{|p{2.8cm}|p{7.3cm}|p{4.2cm}|}
\caption{Yêu cầu phi chức năng và chiến lược đáp ứng}
\label{table:nfr} \\
\hline
\textbf{Yêu cầu} & \textbf{Mô tả} & \textbf{Chiến lược thiết kế} \\ \hline
\endfirsthead
\hline
\textbf{Yêu cầu} & \textbf{Mô tả} & \textbf{Chiến lược thiết kế} \\ \hline
\endhead
Khả năng mở rộng  & Hệ thống có thể mở rộng độc lập từng dịch vụ theo tải
                  & Phân rã 13 service; database per service \\ \hline
Tính sẵn sàng    & Hệ thống tiếp tục hoạt động khi một service gặp sự cố
                  & Circuit Breaker; stateless service; graceful degradation \\ \hline
Tính nhất quán   & Dữ liệu nhất quán cuối cùng giữa các service
                  & Saga Orchestration; Outbox + Idempotent Consumer \\ \hline
Bảo mật          & Xác thực và phân quyền cho mọi API endpoint
                  & Keycloak OAuth2/OIDC; JWT RS256; Gateway Trusted Subsystem \\ \hline
Hiệu năng        & Đáp ứng tải cao trong flash sale; phản hồi nhanh cho catalog
                  & Redis cache; Redis Lua atomic; k6 load test \\ \hline
Khả năng quan sát & Giám sát và truy vết request end-to-end
                  & Prometheus + Grafana; Zipkin; structured logging \\ \hline
Khả năng triển khai & Dễ dàng khởi tạo toàn bộ hệ thống
                  & Docker Compose; healthcheck; Flyway migration tự động \\ \hline
\end{longtable}
}

% =======================================================================
\section{Mô hình hóa yêu cầu --- Biểu đồ Use Case}
\label{section:4.2}

Hệ thống có bốn tác nhân chính tương tác theo các vai trò và mục đích khác nhau. Khách hàng
chưa đăng nhập (Guest) có thể duyệt sản phẩm, tìm kiếm và thao tác với giỏ hàng tạm thời.
Khách hàng đã đăng nhập (Customer) có toàn quyền thực hiện các nghiệp vụ mua sắm: đặt hàng,
thanh toán, tham gia flash sale và đánh giá sản phẩm. Quản trị viên (Admin) quản lý toàn
bộ danh mục, sản phẩm, tồn kho, khuyến mãi, nội dung và chiến dịch flash sale. Hệ thống
bên ngoài (External System) bao gồm VNPAY xử lý thanh toán và SMTP Server gửi email thông
báo. Ngoài ra, các tác vụ nền (Scheduler) tự động xử lý các công việc định kỳ như hết hạn
đặt chỗ tồn kho và timeout thanh toán.

\begin{figure}[!htbp]
    \centering
    \fbox{\parbox{0.85\textwidth}{\centering\vspace{3cm}%
        \texttt{Hinhve/chuong4/usecase.png}%
    \vspace{3cm}}}
    \caption{Biểu đồ Use Case tổng thể hệ thống}
    \label{fig:use-case}
\end{figure}

% =======================================================================
\section{Kiến trúc tổng thể}
\label{section:4.3}

Hệ thống được tổ chức thành bốn tầng kiến trúc rõ ràng, mỗi tầng có trách nhiệm và phạm
vi ảnh hưởng riêng biệt. Tầng Client bao gồm trình duyệt web, ứng dụng di động hoặc các
công cụ kiểm thử API như Postman --- đây là nơi phát sinh yêu cầu vào hệ thống. Tầng Edge
và Spring Cloud chứa API Gateway (cổng vào duy nhất), Eureka Server (đăng ký và khám phá
dịch vụ), Config Server (quản lý cấu hình tập trung) và Keycloak (xác thực và cấp token).
Tầng Business Services gồm 13 dịch vụ nghiệp vụ, mỗi dịch vụ đảm nhiệm một miền chức năng
độc lập và sở hữu cơ sở dữ liệu riêng. Tầng Infrastructure chứa các hệ thống hỗ trợ:
PostgreSQL, Redis, Kafka, Elasticsearch, Mailpit, Prometheus, Grafana và Zipkin.

Giao tiếp giữa các tầng tuân theo nguyên tắc: Client chỉ giao tiếp với tầng Edge (qua API
Gateway); các dịch vụ trong tầng Business giao tiếp với nhau qua REST/OpenFeign hoặc bất
đồng bộ qua Kafka; mọi truy cập đến cơ sở dữ liệu đều thông qua dịch vụ sở hữu dữ liệu đó
--- không có truy cập trực tiếp từ dịch vụ khác.

\begin{figure}[!htbp]
    \centering
    \fbox{\parbox{0.85\textwidth}{\centering\vspace{3cm}%
        \texttt{Hinhve/chuong4/kien-truc-4-tang-he-thong.png}%
    \vspace{3cm}}}
    \caption{Sơ đồ thành phần bốn tầng của hệ thống}
    \label{fig:component-diagram}
\end{figure}

% =======================================================================
\section{Phân rã miền nghiệp vụ và sở hữu dịch vụ}
\label{section:4.4}

Áp dụng nguyên tắc Domain-Driven Design, hệ thống được phân rã thành 13 Bounded Context,
mỗi context tương ứng với một microservice. Bảng~\ref{table:service-ownership} mô tả chi
tiết trách nhiệm, dữ liệu sở hữu và phương thức giao tiếp của từng dịch vụ. Nguyên tắc
cốt lõi là mỗi dịch vụ sở hữu toàn bộ logic nghiệp vụ và dữ liệu trong phạm vi Bounded
Context của mình; mọi truy cập dữ liệu chéo phải thông qua API công khai hoặc sự kiện
được công bố trên Kafka.

{%
\small\setlength{\tabcolsep}{3pt}\renewcommand{\arraystretch}{1.15}%
\begin{longtable}{|p{2.7cm}|p{1.0cm}|p{3.0cm}|p{4.9cm}|p{2.3cm}|}
\caption{Sở hữu dịch vụ và giao tiếp}
\label{table:service-ownership} \\
\hline
\textbf{Dịch vụ} & \textbf{Cổng} & \textbf{Bounded Context} & \textbf{Dữ liệu sở hữu} & \textbf{Giao tiếp} \\ \hline
\endfirsthead
\hline
\textbf{Dịch vụ} & \textbf{Cổng} & \textbf{Bounded Context} & \textbf{Dữ liệu sở hữu} & \textbf{Giao tiếp} \\ \hline
\endhead
identity-service  & 8081 & Định danh \& Đăng ký     & Keycloak users                              & Kafka (pub) \\ \hline
user-service      & 8082 & Hồ sơ người dùng          & PostgreSQL (user\_db)                       & Kafka (sub) \\ \hline
product-service   & 8083 & Sản phẩm \& Danh mục      & PostgreSQL (product\_db), Redis cache       & REST, Kafka (pub) \\ \hline
inventory-service & 8084 & Tồn kho \& Đặt chỗ       & PostgreSQL (inventory\_db)                  & Kafka (pub/sub) \\ \hline
cart-service      & 8085 & Giỏ hàng                  & Redis (cart keys)                           & REST (Feign) \\ \hline
order-service     & 8086 & Đơn hàng (Saga orchestrator) & PostgreSQL (order\_db), Outbox           & REST, Kafka (pub/sub) \\ \hline
payment-service   & 8087 & Thanh toán \& VNPAY       & PostgreSQL (payment\_db), Outbox            & REST, Kafka (pub/sub) \\ \hline
voucher-service   & 8088 & Khuyến mãi                & PostgreSQL (voucher\_db)                    & REST \\ \hline
notification-service & 8089 & Thông báo \& Email    & PostgreSQL (notification\_db)               & Kafka (sub) \\ \hline
review-service    & 8090 & Đánh giá sản phẩm         & PostgreSQL (review\_db)                     & REST (Feign) \\ \hline
search-service    & 8091 & Tìm kiếm toàn văn         & Elasticsearch indices                       & Kafka (sub) \\ \hline
content-service   & 8092 & Nội dung (Banner/Blog)    & PostgreSQL (content\_db)                    & REST \\ \hline
flash-sale-service & 8093 & Flash Sale               & PostgreSQL (flash\_sale\_db), Redis stock   & REST, Kafka (pub/sub) \\ \hline
\end{longtable}
}

\begin{figure}[!htbp]
    \centering
    \fbox{\parbox{0.85\textwidth}{\centering\vspace{3cm}%
        \texttt{Hinhve/chuong4/context-map.png}%
    \vspace{3cm}}}
    \caption{Bản đồ ngữ cảnh (Context Map) 13 Bounded Context}
    \label{fig:context-map}
\end{figure}

% =======================================================================
\section{Thiết kế cơ sở dữ liệu}
\label{section:4.5}

\subsection{PostgreSQL}

Mỗi business service sử dụng PostgreSQL sở hữu một database riêng với schema được quản lý
bởi Flyway migration. Dưới đây là thiết kế cho các database cốt lõi. Các bảng sử dụng UUID
làm khóa chính, đảm bảo tính duy nhất toàn cục và tránh phụ thuộc vào sequence tự tăng của
database --- điều quan trọng trong môi trường phân tán nơi dữ liệu có thể được tạo từ nhiều
service khác nhau.

\subsubsection*{Database đơn hàng --- \texttt{order\_db}}

Đây là database quan trọng nhất, lưu trữ toàn bộ vòng đời của đơn hàng cùng với cơ chế
Outbox và Idempotency. Bảng \texttt{orders} lưu thông tin đơn hàng với các trường chính:
mã đơn, mã khách hàng, phương thức thanh toán, trạng thái, tổng tiền, mã voucher (nếu có)
và thời điểm hết hạn đặt chỗ. Bảng \texttt{order\_items} lưu chi tiết từng mặt hàng trong
đơn: mã sản phẩm, số lượng và đơn giá tại thời điểm đặt. Bảng \texttt{outbox} lưu các sự
kiện cần phát hành lên Kafka, với các trường: mã sự kiện, loại aggregate, mã aggregate, loại
sự kiện, payload JSON, cờ \texttt{processed} và thời điểm xử lý. Bảng
\texttt{processed\_events} lưu mã các sự kiện đã được xử lý, phục vụ cơ chế Idempotent Consumer.

\subsubsection*{Database thanh toán --- \texttt{payment\_db}}

Bảng \texttt{payments} lưu thông tin thanh toán: mã thanh toán, mã đơn hàng, phương thức,
số tiền, trạng thái và thời điểm hết hạn. Bảng \texttt{payment\_outbox} lưu sự kiện thanh
toán ở trạng thái \texttt{PENDING}, số lần thử, lỗi cuối cùng và thời điểm phát hành, đảm bảo
khả năng phát hành sự kiện thanh toán một cách tin cậy.

\subsubsection*{Database sản phẩm --- \texttt{product\_db}}

Bảng \texttt{products} lưu thông tin sản phẩm (tên, mô tả, giá, SKU, trạng thái hoạt động),
liên kết với \texttt{categories} (danh mục phân cấp) và \texttt{brands} (thương hiệu).

\subsubsection*{Database tồn kho --- \texttt{inventory\_db}}

Bảng \texttt{inventory} lưu hai giá trị chính cho từng SKU: \path|quantity| là tổng số lượng
còn trong kho và \path|reserved_quantity| là số lượng đang được đặt chỗ. Số lượng khả dụng
được suy ra bằng \path|quantity - reserved_quantity|. Khi nhận sự kiện
\path|order-created|, inventory-service tăng \path|reserved_quantity|; khi nhận
\path|order-confirmed|, service giảm đồng thời \path|quantity| và
\path|reserved_quantity|; khi nhận \path|order-cancelled|, service giảm
\path|reserved_quantity| để giải phóng đặt chỗ.

\subsubsection*{Các database còn lại}

Các database \texttt{user\_db}, \texttt{voucher\_db}, \texttt{notification\_db},
\texttt{review\_db}, \texttt{content\_db} và \texttt{flash\_sale\_db} lưu trữ dữ liệu
nghiệp vụ tương ứng với mô hình quan hệ chuẩn hóa, mỗi database do một service duy nhất
quản lý.

\begin{figure}[!htbp]
    \centering
    \fbox{\parbox{0.85\textwidth}{\centering\vspace{3cm}%
        \texttt{Hinhve/chuong4/luong-outbox-idempotency.png}%
    \vspace{3cm}}}
    \caption{Luồng hoạt động của Transactional Outbox và Idempotent Consumer}
    \label{fig:outbox-idempotency-flow}
\end{figure}

\begin{figure}[!htbp]
    \centering
    \fbox{\parbox{0.85\textwidth}{\centering\vspace{3cm}%
        \texttt{Hinhve/chuong4/erd-main.png}%
    \vspace{3cm}}}
    \caption{Mô hình ERD các database chính: order\_db, payment\_db, product\_db, inventory\_db}
    \label{fig:erd-main}
\end{figure}

\subsection{Elasticsearch}

Search-service sử dụng Elasticsearch để lưu trữ và truy vấn thông tin sản phẩm phục vụ tìm
kiếm toàn văn. Mỗi sản phẩm được lưu dưới dạng một document trong index \texttt{products}
với cấu trúc được tối ưu cho tìm kiếm: trường \texttt{name} và \texttt{description} được
phân tích toàn văn (text analyzer với stemming tiếng Việt nếu có), trường \texttt{category}
và \texttt{brand} được lưu dưới dạng keyword để lọc chính xác, trường \texttt{price} lưu
dưới dạng số để lọc theo khoảng giá. Document được đồng bộ từ PostgreSQL của product-service
thông qua Kafka event mỗi khi có thay đổi, đảm bảo eventual consistency giữa hai mô hình.

\subsection{Redis}

Redis được sử dụng cho ba mục đích với cấu trúc khóa được thiết kế rõ ràng. Đối với giỏ
hàng, khóa có định dạng \path|cart:{userId}| hoặc \path|cart:{sessionId}| lưu trữ
danh sách sản phẩm dưới dạng Hash, với TTL 7 ngày cho người dùng đã đăng nhập và 30 phút
cho khách vãng lai. Đối với cache sản phẩm, khóa \path|product:{productId}| lưu thông
tin sản phẩm dưới dạng JSON string với TTL 30 phút, được tự động invalidate khi sản phẩm
được cập nhật. Đối với flash sale, khóa \path|flash-sale:{campaignId}:stock| lưu số
lượng slot còn lại dưới dạng integer, được khởi tạo khi chiến dịch được kích hoạt và tự động
hết hạn theo TTL bằng với thời gian còn lại của chiến dịch; khóa
\path|flash-sale:{campaignId}:buyers| lưu tập hợp (Set) mã người dùng đã mua thành công
để chống trùng lặp.

\begin{figure}[!htbp]
    \centering
    \fbox{\parbox{0.85\textwidth}{\centering\vspace{3cm}%
        \texttt{Hinhve/chuong4/redis.png}%
    \vspace{3cm}}}
    \caption{Mô hình khóa Redis: Cart, Cache, Flash Sale}
    \label{fig:redis-keys}
\end{figure}

% =======================================================================
\section{Thiết kế API và giao tiếp đồng bộ}
\label{section:4.6}

\textbf{Định tuyến API Gateway.} API Gateway định tuyến yêu cầu từ client đến các dịch vụ
backend dựa trên đường dẫn. Ngoài
chức năng định tuyến, Gateway còn áp dụng các chính sách bảo mật và giới hạn tốc độ khác
nhau cho từng nhóm endpoint. Bảng~\ref{table:gateway-routes} tổng hợp các route chính.

{%
\small\setlength{\tabcolsep}{3pt}\renewcommand{\arraystretch}{1.15}%
\begin{longtable}{|p{3.8cm}|p{3.0cm}|p{4.7cm}|p{2.6cm}|}
\caption{Định tuyến API Gateway}
\label{table:gateway-routes} \\
\hline
\textbf{Đường dẫn} & \textbf{Dịch vụ đích} & \textbf{Xác thực} & \textbf{Rate Limit} \\ \hline
\endfirsthead
\hline
\textbf{Đường dẫn} & \textbf{Dịch vụ đích} & \textbf{Xác thực} & \textbf{Rate Limit} \\ \hline
\endhead
\path|/api/auth/**|              & identity-service     & Không (public)              & 10 req/s, burst 20 \\ \hline
\path|/api/products/**|          & product-service      & Tùy chọn (GET public)       & Không giới hạn \\ \hline
\path|/api/search/**|            & search-service       & Tùy chọn                   & Không giới hạn \\ \hline
\path|/api/cart/**|              & cart-service         & Tùy chọn (guest session)    & Không giới hạn \\ \hline
\path|/api/users/**|             & user-service         & Bắt buộc                    & Không giới hạn \\ \hline
\path|/api/orders/**|            & order-service        & Bắt buộc                    & 10 req/s, burst 20 \\ \hline
\path|/api/payments/**|          & payment-service      & Bắt buộc (trừ VNPAY IPN)   & Không giới hạn \\ \hline
\path|/api/vouchers/**|          & voucher-service      & Bắt buộc (admin cho CRUD)  & Không giới hạn \\ \hline
\path|/api/reviews/**|           & review-service       & Bắt buộc (tạo), tùy chọn (xem) & Không giới hạn \\ \hline
\path|/api/flash-sales/*/purchase| & flash-sale-service & Bắt buộc                   & 3 req/s, burst 5 \\ \hline
\path|/api/flash-sales/**|       & flash-sale-service   & Bắt buộc (admin cho CRUD)  & Không giới hạn \\ \hline
\path|/api/content/**|           & content-service      & Tùy chọn (GET public)       & Không giới hạn \\ \hline
\path|/api/notifications/**|     & notification-service & Bắt buộc                    & Không giới hạn \\ \hline
\end{longtable}
}

\textbf{Giao tiếp Feign giữa các dịch vụ.} Các tương tác đồng bộ giữa các dịch vụ được thực
hiện thông qua OpenFeign client. Cart-service
gọi product-service để xác minh tính hợp lệ của sản phẩm khi thêm vào giỏ hàng. Order-service
--- với vai trò Saga orchestrator --- gọi product-service để kiểm tra giá và trạng thái sản
phẩm, gọi voucher-service để xác minh và đặt chỗ voucher khi tạo đơn, và gọi cart-service để
xóa giỏ hàng sau khi đơn được xác nhận. Payment-service gọi order-service để lấy payment
context; review-service gọi order-service để kiểm tra người dùng đã có đơn hàng đã xác nhận
chứa sản phẩm cần đánh giá hay chưa; flash-sale-service gọi order-service để đếm số đơn phục
vụ reconciliation. Việc đặt chỗ tồn kho không dùng Feign mà đi qua sự kiện Kafka
\texttt{order-created} trong luồng Saga. Các Feign client được cấu hình timeout 2--3 giây và
retry 2--3 lần (tùy client), kết hợp với Circuit Breaker để ngăn cascading failure.

\begin{figure}[!htbp]
    \centering
    \fbox{\parbox{0.85\textwidth}{\centering\vspace{3cm}%
        \texttt{Hinhve/chuong4/feign-dependencies.png}%
    \vspace{3cm}}}
    \caption{Sơ đồ phụ thuộc REST/Feign giữa các dịch vụ}
    \label{fig:feign-dependencies}
\end{figure}

% =======================================================================
\section{Thiết kế luồng sự kiện --- Kafka Topics}
\label{section:4.7}

Giao tiếp bất đồng bộ giữa các dịch vụ được thực hiện thông qua 13 Kafka topic, được tổ
chức thành 11 nhóm nghiệp vụ. Mỗi topic có producer và consumer được xác định rõ ràng, cùng
với cấu trúc payload cụ thể. Kafka được cấu hình với cơ chế tự động tạo topic
(\texttt{auto.create.topics.enable=true}) cho hầu hết các topic, ngoại trừ
\texttt{flash-sale-order-requested} được tạo thủ công với 3 partition để kiểm soát thứ tự
xử lý.

Bảng~\ref{table:kafka-topics-design} mô tả toàn bộ các topic, luồng sự kiện và vai trò của từng
dịch vụ tham gia.

{%
\small\setlength{\tabcolsep}{3pt}\renewcommand{\arraystretch}{1.15}%
\begin{longtable}{|p{3.3cm}|p{2.6cm}|p{3.6cm}|p{4.6cm}|}
\caption{Bản đồ Kafka Topics}
\label{table:kafka-topics-design} \\
\hline
\textbf{Topic} & \textbf{Producer} & \textbf{Consumer} & \textbf{Mục đích} \\ \hline
\endfirsthead
\hline
\textbf{Topic} & \textbf{Producer} & \textbf{Consumer} & \textbf{Mục đích} \\ \hline
\endhead
\path|user-registered|           & identity-service  & user-service               & Tạo hồ sơ người dùng sau đăng ký \\ \hline
\path|product-created|           & product-service   & search-service             & Index sản phẩm mới vào Elasticsearch \\ \hline
\path|product-updated|           & product-service   & search-service             & Cập nhật index khi sửa sản phẩm \\ \hline
\path|product-deleted|           & product-service   & search-service             & Xóa index khi xóa sản phẩm \\ \hline
\path|order-created|             & order-service     & inventory-service          & Yêu cầu đặt chỗ tồn kho \\ \hline
\path|inventory-updated|         & inventory-service & order-service              & Xác nhận đặt chỗ thành công \\ \hline
\path|inventory-failed|          & inventory-service & order-service              & Thông báo không đủ tồn kho \\ \hline
\path|payment-requested|         & order-service     & payment-service            & Yêu cầu xử lý thanh toán COD \\ \hline
\path|payment-success|           & payment-service   & order-service              & Thanh toán thành công, xác nhận đơn \\ \hline
\path|payment-failed|            & payment-service   & order-service              & Thanh toán thất bại, hủy đơn \\ \hline
\path|order-confirmed|           & order-service     & inventory-service, notification-service & Xác nhận tồn kho, gửi email \\ \hline
\path|order-cancelled|           & order-service     & inventory-service, notification-service & Hoàn tồn kho, gửi email (voucher hoàn qua Feign) \\ \hline
\path|flash-sale-order-requested| & flash-sale-service & order-service             & Tạo đơn hàng flash sale \\ \hline
\end{longtable}
}

\begin{figure}[!htbp]
    \centering
    \fbox{\parbox{0.85\textwidth}{\centering\vspace{3cm}%
        \texttt{Hinhve/chuong4/kafka-flow.png}%
    \vspace{3cm}}}
    \caption{Sơ đồ luồng sự kiện Kafka giữa các dịch vụ}
    \label{fig:kafka-flow}
\end{figure}

% =======================================================================
\section{Thiết kế bảo mật}
\label{section:4.8}

Hệ thống áp dụng mô hình bảo mật tập trung với API Gateway là điểm thực thi duy nhất. Luồng
xác thực diễn ra qua bốn bước. Bước thứ nhất, client gửi thông tin đăng nhập đến Keycloak
(qua identity-service hoặc trực tiếp) và nhận được bộ token JWT. Bước thứ hai, client gửi
kèm Access Token trong header \texttt{Authorization: Bearer <token>} cho mọi request đến
API Gateway. Bước thứ ba, Gateway xác minh chữ ký JWT sử dụng public key lấy từ endpoint
JWKS của Keycloak, kiểm tra thời hạn và các claims cần thiết. Bước thứ tư, sau khi xác minh
thành công, Gateway trích xuất thông tin định danh từ JWT (userId, roles, email) và chèn
vào ba header \texttt{X-User-Id}, \texttt{X-User-Roles}, \texttt{X-User-Email} trước khi
chuyển tiếp request đến dịch vụ backend.

Các dịch vụ backend không thực hiện lại việc xác minh JWT mà hoàn toàn tin tưởng vào các
header do Gateway cung cấp (Trusted Subsystem Pattern). Điều kiện để mô hình này an toàn
là các dịch vụ backend chỉ lắng nghe trên mạng nội bộ của Docker, không có cổng nào được
mở ra bên ngoài. Phân quyền được thực hiện dựa trên vai trò: các route quản trị như
\path|/api/orders/admin/**|, \path|/api/users/admin/**|, \path|/api/inventory/admin/**|,
\path|/api/reviews/admin/**| và các mutation product/voucher/flash-sale/content yêu cầu
\path|ROLE_ADMIN|; các endpoint nghiệp vụ người dùng yêu cầu token hợp lệ hoặc được mở
public tùy chức năng.

Các endpoint đặc biệt nhạy cảm được bảo vệ bổ sung bằng Redis token bucket rate limiter.
Endpoint đặt hàng (\texttt{POST /api/orders}) được giới hạn 10 request/giây với khả năng
burst 20 request để ngăn chặn đặt hàng tự động. Endpoint mua flash sale
(\texttt{POST /api/flash-sales/*/purchase}) được giới hạn 3 request/giây với burst 5
request, đảm bảo công bằng giữa những người tham gia.

\begin{figure}[!htbp]
    \centering
    \fbox{\parbox{0.85\textwidth}{\centering\vspace{3cm}%
        \texttt{Hinhve/chuong4/keycloak.png}%
    \vspace{3cm}}}
    \caption{Luồng xác thực và phân quyền qua API Gateway}
    \label{fig:security-flow}
\end{figure}

% =======================================================================
\section{Thiết kế máy trạng thái}
\label{section:4.9}

\textbf{Máy trạng thái đơn hàng.} Đơn hàng trong hệ thống trải qua một vòng đời được kiểm
soát chặt chẽ với bốn trạng thái chính. Trạng thái \texttt{PENDING} là trạng thái khởi tạo, ngay sau khi order-service tạo
đơn hàng và ghi sự kiện \texttt{order-created} vào outbox. Trạng thái
\texttt{STOCK\_RESERVED} đạt được khi inventory-service xác nhận đã đặt chỗ tồn kho thành
công. Từ đây, luồng xử lý rẽ nhánh tùy theo phương thức thanh toán: với COD, đơn hàng được
xác nhận ngay lập tức; với VNPAY, đơn hàng chờ khách hàng thanh toán trong vòng 30 phút.
Trạng thái \texttt{CONFIRMED} là trạng thái kết thúc thành công, đạt được khi thanh toán
COD được xác nhận hoặc VNPAY callback thành công. Trạng thái \texttt{CANCELLED} là trạng
thái kết thúc thất bại, xảy ra khi tồn kho không đủ, thanh toán thất bại hoặc hết thời
gian chờ thanh toán.

\begin{figure}[!htbp]
    \centering
    \fbox{\parbox{0.85\textwidth}{\centering\vspace{3cm}%
        \texttt{Hinhve/chuong4/order-service-state-machine.png}%
    \vspace{3cm}}}
    \caption{Sơ đồ máy trạng thái đơn hàng}
    \label{fig:order-state-machine}
\end{figure}

\textbf{Máy trạng thái thanh toán.} Thanh toán VNPAY có bốn trạng thái. \texttt{PENDING} --- URL thanh toán đã được tạo, khách
hàng đang trong quá trình thanh toán. \texttt{COMPLETED} --- VNPAY xác nhận thanh toán
thành công qua IPN callback hoặc Return URL. \texttt{FAILED} --- thanh toán thất bại do
khách hàng hủy, thẻ không hợp lệ hoặc số dư không đủ. \texttt{TIMEOUT} --- quá thời gian
chờ 30 phút, được phát hiện và xử lý bởi PaymentTimeoutScheduler. Một ràng buộc quan trọng
là mỗi đơn hàng chỉ có thể có tối đa một thanh toán ở trạng thái \texttt{PENDING} tại mỗi
thời điểm, được đảm bảo bằng partial unique index trên \texttt{order\_id} khi
\texttt{status = 'PENDING'}.

\textbf{Máy trạng thái chiến dịch Flash Sale.} Chiến dịch flash sale vận hành qua bốn trạng thái. \texttt{SCHEDULED} --- chiến dịch đã được
tạo và lên lịch, đang chờ đến thời điểm bắt đầu. \texttt{ACTIVE} --- chiến dịch đang diễn
ra, người dùng có thể tham gia mua hàng; ở trạng thái này, CampaignScheduler đã khởi tạo
Redis stock key và TTL. \texttt{ENDED} --- chiến dịch đã kết thúc theo thời gian, không nhận
thêm yêu cầu mua mới. \texttt{CANCELLED} --- chiến dịch bị quản trị viên hủy trước hoặc trong
thời gian diễn ra. Quá trình chuyển trạng thái từ \texttt{SCHEDULED} sang \texttt{ACTIVE}
được thực hiện tự động bởi CampaignScheduler (chu kỳ 5 giây) dựa trên thời gian hệ thống.
Việc chuyển sang \texttt{ENDED} cũng được CampaignScheduler xử lý khi thời gian kết thúc đã
qua; khi hủy chiến dịch, service xóa các Redis key liên quan để không tiếp nhận mua mới.

\begin{figure}[!htbp]
    \centering
    \fbox{\parbox{0.85\textwidth}{\centering\vspace{3cm}%
        \texttt{Hinhve/chuong4/flashsale-state-machine.png}%
    \vspace{3cm}}}
    \caption{Sơ đồ máy trạng thái thanh toán và chiến dịch Flash Sale}
    \label{fig:payment-flashsale-state-machine}
\end{figure}

% =======================================================================
\section{Thiết kế luồng nghiệp vụ chính}
\label{section:4.10}

\subsection{Luồng thao tác giỏ hàng}

Giỏ hàng trong hệ thống được thiết kế để hỗ trợ cả người dùng đã đăng nhập và khách vãng lai
thông qua cơ chế session. Khi người dùng thêm sản phẩm vào giỏ hàng qua
\path|POST /api/cart/items|, cart-service gọi product-service qua OpenFeign để xác minh
sản phẩm tồn tại và đang hoạt động, sau đó lưu thông tin vào Redis với key được tạo từ
userId (nếu đã đăng nhập) hoặc sessionId (nếu là khách). Khi người dùng đăng nhập, session
cart được tự động hợp nhất (merge) vào user cart thông qua API
\path|POST /api/cart/merge|. Dữ liệu giỏ hàng được đặt TTL 7 ngày cho người dùng đã đăng
nhập và 30 phút cho khách vãng lai, tự động được dọn dẹp khi hết hạn.

\begin{figure}[!htbp]
    \centering
    \fbox{\parbox{0.85\textwidth}{\centering\vspace{3cm}%
        \texttt{Hinhve/chuong4/luong-gio-hang.png}%
    \vspace{3cm}}}
    \caption{Biểu đồ tuần tự luồng thao tác giỏ hàng}
    \label{fig:sequence-cart}
\end{figure}

\subsection{Luồng đăng ký người dùng}

Khi người dùng gửi yêu cầu đăng ký qua \path|POST /api/auth/register|, identity-service
tiếp nhận và gọi Keycloak Admin API để tạo tài khoản mới trong realm \texttt{ecommerce}.
Sau khi tạo thành công, identity-service phát hành sự kiện \texttt{user-registered} lên
Kafka, chứa thông tin: userId, email, username và timestamp. User-service tiêu thụ sự kiện
này và tạo bản ghi User profile tương ứng trong \texttt{user\_db}, bao gồm các trường mặc
định như họ tên (từ email), trạng thái hoạt động và thời điểm tạo. Luồng này minh họa mẫu
giao tiếp bất đồng bộ: identity-service không cần biết user-service tồn tại, và user-service
có thể được triển khai hoặc thay đổi độc lập.

\begin{figure}[!htbp]
    \centering
    \fbox{\parbox{0.85\textwidth}{\centering\vspace{3cm}%
        \texttt{Hinhve/chuong4/dangki-dangnhap.png}%
    \vspace{3cm}}}
    \caption{Biểu đồ tuần tự luồng đăng kí và đăng nhập người dùng}
    \label{fig:sequence-register-login}
\end{figure}

\subsection{Luồng đặt hàng COD --- Saga Orchestration}

Đây là luồng phức tạp nhất, thể hiện đầy đủ Saga Orchestration Pattern với order-service
đóng vai trò orchestrator. Luồng diễn ra qua các bước sau:

\textbf{Bước 1 --- Tạo đơn hàng.} Khách hàng gửi yêu cầu \path|POST /api/orders| với
payload chứa danh sách sản phẩm (productId, quantity), phương thức thanh toán COD, địa
chỉ giao hàng và mã voucher (nếu có). Order-service thực hiện hai lời gọi Feign đồng bộ:
gọi product-service để xác minh sản phẩm tồn tại và đang hoạt động, lấy đơn giá hiện tại;
và gọi voucher-service để xác minh và đặt chỗ voucher (nếu có). Việc kiểm tra và đặt chỗ
tồn kho không thực hiện đồng bộ ở bước này mà được xử lý bất đồng bộ ở Bước 2 qua sự kiện
Kafka. Nếu các lời gọi đều hợp lệ, order-service
tạo bản ghi đơn hàng với trạng thái \texttt{PENDING} và đồng thời ghi sự kiện
\texttt{order-created} vào bảng \texttt{outbox} trong cùng một giao dịch.

\textbf{Bước 2 --- Đặt chỗ tồn kho.} OutboxPoller (chu kỳ 1 giây) của order-service phát
hành sự kiện \texttt{order-created} lên Kafka. Inventory-service tiêu thụ sự kiện, thực
hiện đặt chỗ cho từng sản phẩm trong đơn hàng bằng cách tăng \texttt{reserved\_quantity}
trong khi \texttt{quantity} chưa đổi. Nếu tất cả sản phẩm đều đủ số lượng, inventory-service phát
hành \texttt{inventory-updated} lên Kafka. Nếu bất kỳ sản phẩm nào không đủ, phát hành
\texttt{inventory-failed}.

\begin{figure}[!htbp]
    \centering
    \fbox{\parbox{0.85\textwidth}{\centering\vspace{3cm}%
        \texttt{Hinhve/chuong4/luong-taodon-reservetonkho.png}%
    \vspace{3cm}}}
    \caption{Biểu đồ tuần tự Bước 1--2 Saga: Tạo đơn hàng và đặt chỗ tồn kho qua Transactional Outbox}
    \label{fig:sequence-create-reserve}
\end{figure}

\textbf{Bước 3 --- Xác nhận đơn COD.} Order-service tiêu thụ \path|inventory-updated|,
chuyển trạng thái đơn hàng sang \texttt{STOCK\_RESERVED}. Với phương thức COD, quá trình
thanh toán được xử lý tự động: order-service phát hành \texttt{order-confirmed} ngay lập
tức. Inventory-service tiêu thụ \path|order-confirmed|, thực hiện confirm reservation
(trừ hẳn số lượng tồn kho). Notification-service tiêu thụ \texttt{order-confirmed}, gửi
email xác nhận đơn hàng đến khách hàng. Đơn hàng đạt trạng thái cuối cùng \texttt{CONFIRMED}.

\textbf{Bước 4 --- Xử lý thất bại.} Nếu inventory-service phát hành \texttt{inventory-failed},
order-service chuyển trạng thái đơn hàng sang \texttt{CANCELLED} với lý do không đủ tồn
kho, phát hành \texttt{order-cancelled}. Inventory-service và voucher-service tiêu thụ
\texttt{order-cancelled} để hoàn trả đặt chỗ và voucher. Notification-service gửi email
thông báo hủy đơn.

\begin{figure}[!htbp]
    \centering
    \fbox{\parbox{0.85\textwidth}{\centering\vspace{3cm}%
        \texttt{Hinhve/chuong4/luong-COD-payment-success.png}%
    \vspace{3cm}}}
    \caption{Biểu đồ tuần tự luồng đặt hàng COD --- Saga Orchestration}
    \label{fig:sequence-cod}
\end{figure}

\subsection{Luồng thông báo email}

Notification-service đóng vai trò trung tâm trong việc gửi thông báo đến khách hàng, hoạt
động hoàn toàn theo cơ chế event-driven. Dịch vụ lắng nghe hai sự kiện chính từ Kafka:
\texttt{order-confirmed} để gửi email xác nhận đơn hàng kèm chi tiết sản phẩm và tổng tiền;
và \texttt{order-cancelled} để gửi email thông báo hủy đơn kèm lý do. Mỗi email được
tạo từ template với các trường động được điền từ payload của sự kiện, sau đó gửi qua SMTP
server --- trong môi trường phát triển là Mailpit, trong môi trường production là dịch vụ
email thật.

\begin{figure}[!htbp]
    \centering
    \fbox{\parbox{0.85\textwidth}{\centering\vspace{3cm}%
        \texttt{Hinhve/chuong4/luong-thongbao-email.png}%
    \vspace{3cm}}}
    \caption{Biểu đồ tuần tự luồng gửi thông báo email qua notification-service}
    \label{fig:sequence-email}
\end{figure}

\subsection{Luồng đặt hàng VNPAY}

Luồng VNPAY mở rộng luồng COD với các bước bổ sung cho thanh toán trực tuyến. Sau khi đơn
hàng đạt trạng thái \path|STOCK_RESERVED|, order-service thiết lập
\path|reservation_expired_at = NOW() + 30| phút. Khách hàng gọi
\path|POST /api/payments/vnpay/create| để tạo URL thanh toán. Payment-service sinh URL
với chữ ký HMAC-SHA512 và trả về cho client. Khách hàng được chuyển hướng đến trang thanh
toán VNPAY.

Khi VNPAY gọi IPN callback với kết quả thanh toán, payment-service xác minh chữ ký và cập
nhật trạng thái. Nếu thành công, payment-service phát hành \texttt{payment-success} lên
Kafka; order-service tiêu thụ và chuyển đơn hàng sang \texttt{CONFIRMED}. Nếu thất bại,
phát hành \texttt{payment-failed} và đơn hàng bị hủy với cơ chế hoàn trả tương tự. Nếu sau
30 phút không có callback, ReservationExpiryScheduler (order-service) và
PaymentTimeoutScheduler (payment-service) cùng hoạt động từ hai phía để đảm bảo đơn hàng
và thanh toán đều được chuyển sang trạng thái kết thúc phù hợp.

\begin{figure}[!htbp]
    \centering
    \fbox{\parbox{0.85\textwidth}{\centering\vspace{3cm}%
        \texttt{Hinhve/chuong4/luong-thanhtoan-thatbai.png}%
    \vspace{3cm}}}
    \caption{Biểu đồ tuần tự luồng xử lý thanh toán thất bại và hoàn trả tài nguyên}
    \label{fig:sequence-payment-fail}
\end{figure}

\begin{figure}[!htbp]
    \centering
    \fbox{\parbox{0.85\textwidth}{\centering\vspace{3cm}%
        \texttt{Hinhve/chuong4/luong-VNPay-callback-thanhcong.png}%
    \vspace{3cm}}}
    \caption{Biểu đồ tuần tự luồng đặt hàng VNPAY với xử lý IPN callback}
    \label{fig:sequence-vnpay}
\end{figure}

\subsection{Luồng Flash Sale}

Luồng flash sale được thiết kế đặc biệt để xử lý tình huống đồng thời cao với nhiều lớp bảo vệ.
Khi khách hàng gọi \path|POST /api/flash-sales/{campaignId}/purchase|, flash-sale-service
thực hiện tuần tự các bước kiểm tra. Đầu tiên, kiểm tra thời gian hiệu lực của chiến dịch
(phải ở trạng thái ACTIVE và trong khoảng thời gian cho phép). Tiếp theo, thực thi Redis Lua
script nguyên tử để kiểm tra buyer set và giảm số lượng slot: script kiểm tra khóa
\path|flash-sale:{campaignId}:buyers| để chống mua trùng, sau đó kiểm tra khóa
\path|flash-sale:{campaignId}:stock|; nếu stock lớn hơn 0 thì giảm 1 và thêm user vào
buyer set, ngược lại trả về mã thất bại. Nếu Redis trả về thất bại, service ném lỗi nghiệp
vụ (HTTP 400) với thông điệp tương ứng --- sản phẩm đã bán hết hoặc người dùng đã mua trước đó.

Nếu Redis trả về thành công, flash-sale-service phát hành sự kiện
\path|flash-sale-order-requested| lên Kafka rồi tăng \texttt{soldCount} trong database (nếu
bước tăng \texttt{soldCount} lỗi sau khi Kafka đã ACK, ReconciliationScheduler sẽ bù lại sau).
Order-service tạo đơn hàng flash sale với giá đã giảm theo chiến dịch và xử lý tương tự luồng
COD. Nếu bước tạo đơn thất bại, flash-sale-service dùng compensation script để hoàn lại slot
Redis và loại người dùng khỏi buyer set. ReconciliationScheduler định kỳ đối chiếu
\texttt{soldCount} với số đơn flash-sale thực tế trong order-service để phát hiện và bù trừ
slot mồ côi.

CampaignScheduler (chu kỳ 5 giây, với Redis distributed lock) đảm nhiệm ba nhiệm vụ nền:
kích hoạt chiến dịch đến hạn (chuyển \texttt{SCHEDULED} sang \texttt{ACTIVE}, seed Redis
stock), khôi phục khóa Redis bị mất (đọc \texttt{soldCount} từ database để tính lại stock
còn lại), và kết thúc chiến dịch hết hạn (chuyển \texttt{ACTIVE} sang \texttt{ENDED}).

\begin{figure}[!htbp]
    \centering
    \fbox{\parbox{0.85\textwidth}{\centering\vspace{3cm}%
        \texttt{Hinhve/chuong4/luong-flash-sale.png}%
    \vspace{3cm}}}
    \caption{Biểu đồ tuần tự luồng mua hàng Flash Sale với cơ chế ba lớp bảo vệ}
    \label{fig:sequence-flashsale}
\end{figure}

\subsection{Luồng tìm kiếm sản phẩm (CQRS-lite)}

Luồng tìm kiếm minh họa áp dụng CQRS-lite với hai đường dữ liệu độc lập. Đường ghi
(Command): khi quản trị viên tạo hoặc cập nhật sản phẩm qua product-service, dữ liệu được
lưu vào PostgreSQL (\texttt{product\_db}). Đồng thời, product-service phát hành sự kiện
\texttt{product-created} hoặc \texttt{product-updated} trực tiếp lên Kafka. Search-service
tiêu thụ sự kiện và cập nhật document tương ứng trong Elasticsearch index. Đường đọc
(Query): khi khách hàng tìm kiếm sản phẩm, request được gửi trực tiếp đến search-service,
truy vấn Elasticsearch và trả về kết quả đã được tối ưu cho tìm kiếm toàn văn, không thông
qua PostgreSQL. Hai đường dữ liệu hoạt động độc lập, chấp nhận độ trễ đồng bộ nhỏ (thường
dưới 1 giây), đảm bảo eventual consistency.

\begin{figure}[!htbp]
    \centering
    \fbox{\parbox{0.85\textwidth}{\centering\vspace{3cm}%
        \texttt{Hinhve/chuong4/luong-tim-kiem-sp.png}%
    \vspace{3cm}}}
    \caption{Biểu đồ tuần tự luồng tìm kiếm sản phẩm theo mô hình CQRS-lite}
    \label{fig:sequence-cqrs}
\end{figure}

\subsection{Luồng đánh giá sản phẩm}

Luồng đánh giá sản phẩm minh họa cơ chế xác minh quyền đánh giá dựa trên lịch sử mua hàng.
Khi khách hàng gửi yêu cầu \path|POST /api/reviews|, review-service gọi order-service
qua OpenFeign để kiểm tra người dùng đã có đơn hàng ở trạng thái \texttt{CONFIRMED} chứa
sản phẩm cần đánh giá hay chưa. Chỉ khi điều kiện này được thỏa mãn, đánh giá mới được
tạo và lưu vào \texttt{review\_db}. Cơ chế này ngăn chặn đánh giá giả mạo từ người dùng
chưa từng mua sản phẩm, đảm bảo tính xác thực của hệ thống đánh giá.

\begin{figure}[!htbp]
    \centering
    \fbox{\parbox{0.85\textwidth}{\centering\vspace{3cm}%
        \texttt{Hinhve/chuong4/luong-review-sp.png}%
    \vspace{3cm}}}
    \caption{Biểu đồ tuần tự luồng đánh giá sản phẩm}
    \label{fig:sequence-review}
\end{figure}

% =======================================================================
\section{Thiết kế tác vụ nền --- Scheduled Jobs}
\label{section:4.11}

Hệ thống có sáu tác vụ nền (scheduled jobs) chạy định kỳ để xử lý các tình huống mà cơ chế
event-driven không bao phủ được, đặc biệt là các tình huống liên quan đến timeout và phục
hồi. Bảng~\ref{table:scheduled-jobs} tổng hợp toàn bộ các tác vụ.

{%
\small\setlength{\tabcolsep}{3pt}\renewcommand{\arraystretch}{1.15}%
\begin{longtable}{|p{3.9cm}|p{2.7cm}|p{1.9cm}|p{5.6cm}|}
\caption{Tác vụ nền trong hệ thống}
\label{table:scheduled-jobs} \\
\hline
\textbf{Tác vụ} & \textbf{Dịch vụ} & \textbf{Chu kỳ} & \textbf{Nhiệm vụ} \\ \hline
\endfirsthead
\hline
\textbf{Tác vụ} & \textbf{Dịch vụ} & \textbf{Chu kỳ} & \textbf{Nhiệm vụ} \\ \hline
\endhead
\makecell[l]{ReservationExpiry\\Scheduler} & order-service   & 60 giây  & Hủy đơn VNPAY quá hạn 30 phút, giải phóng voucher \\ \hline
\makecell[l]{PaymentTimeout\\Scheduler}    & payment-service & 60 giây  & Chuyển payment PENDING quá hạn sang TIMEOUT \\ \hline
OutboxPoller (Order)                        & order-service   & 1 giây   & Phát hành sự kiện chưa gửi từ bảng outbox \\ \hline
OutboxPoller (Payment)                      & payment-service & 1 giây   & Phát hành sự kiện chưa gửi từ bảng payment\_outbox \\ \hline
\makecell[l]{Campaign\\Scheduler}          & flash-sale-service & 5 giây & Kích hoạt campaign, khôi phục Redis key, kết thúc campaign hết hạn \\ \hline
\makecell[l]{Reconciliation\\Scheduler}    & flash-sale-service & 300 giây & Đối chiếu soldCount với số đơn flash-sale thực tế và bù slot mồ côi \\ \hline
\end{longtable}
}

Hai scheduler xử lý hết hạn reservation và timeout payment hoạt động song song để xử lý
cùng một tình huống từ hai phía, tạo cơ chế dự phòng chéo: nếu một scheduler gặp sự cố hoặc
bỏ sót, scheduler còn lại vẫn đảm bảo tài nguyên (tồn kho, voucher) được giải phóng.
OutboxPoller với chu kỳ 1 giây đảm bảo độ trễ phát hành sự kiện thấp trong điều kiện
bình thường, đồng thời phục hồi các sự kiện bị bỏ lỡ sau khi service khởi động lại.
CampaignScheduler sử dụng Redis distributed lock (SETNX với TTL) để đảm bảo chỉ một
instance của flash-sale-service thực thi tại mỗi thời điểm, tránh xung đột trong môi
trường nhiều replica.

\end{document}
